\section{一八七五至一九一二：改革、教育、立宪与省份社会}
\label{sec:late-qing-reform-constitution-provincial-global}

同治末至帝制终结的改革，不是由北京发布并在各省均匀展开的一张清单。新疆、台湾、东北、直隶、江南、两湖、两广、西南及口岸城市面对不同战争、财政、交通和社会组织；学堂、警察、法院、新军、商会、咨议局与报刊也以不同速度进入地方。实录、军机档、奏折、章程和官报能说明制度被宣布，学校名册、预算、地方档案、报刊、契约、教会材料与外文外交档案则只能校验某些地点的执行。创设、经费、实际运作和居民经验不可混为一事。

\subsection{边疆战争、设省与地方治理}

一八七五年光绪即位后的西北军政、左宗棠收复新疆、伊犁谈判和一八八四年新疆设省，显示军事胜利、外交文本与日常行政之间存在长期过渡。清廷档案可以确认设官、筹饷和条约，不能证明绿洲、牧区、城镇和山地居民拥有同一种司法、赋役或教育经验。维吾尔、蒙古、回、满、汉等称谓在军政、宗教和地方材料中具有不同语境，不应被当作固定、内部无差异的现代族群。

一八八五年台湾建省及铁路、电报、矿务、海防等项目，亦不能以开办日期替代地方社会的变化。预算拨付、土地征用、工人招募、线路使用和不同社群的承受需分别核验。省级建制、港口和内山之间并非同一行政空间；将“现代化”写成自动向内陆扩散，会遮蔽道路、财政和地方权力的限制。

\subsection{教育、性别与职业的新资格}

一八九五年以后新式学堂、翻译、留学、报刊与实业教育更为突出，一九〇五年废科举重写了入仕与读书的制度语言。章程和上谕可以确证改革目标，学校是否有稳定经费、教师、课本、女生与学生出路则因地方而异。科举废除不意味着乡村私塾立即消失，也不表示女性、贫民或边疆居民获得同等进入新学的机会。

女学、妇女报刊、慈善、工厂和城市职业的材料使部分女性的教育与公共活动可见。它们主要来自城市、有文字能力者或改革倡导者，不能代表所有家庭。婚姻、继承、嫁资、劳动与迁徙仍受家庭、法律、市场和地方习惯共同影响；不应以“新女性”或“传统妇女”二分复杂而不均衡的经验。

\subsection{法律、警政与地方公共性}

一九〇六年预备立宪、一九〇七年修律与警政扩展、一九〇八年《钦定宪法大纲》，说明朝廷尝试重写国家制度。律例修订、警察局和审判机构的设立，不能证明所有县份已拥有同样的人员、经费和程序。旧有州县、书吏、绅商、宗族和会馆仍会参与调解、治安和征收；新制度的名称与实际权力之间需要个案材料检验。

商会、地方自治讨论、咨议局和报刊形成新的公共表述空间，也让地方精英、商人、学生和官员以不同方式争论税收、铁路、教育和宪政。报刊的公开意见并非人口抽样，议员与商会成员也不等于全体居民。研究省份政治须同时看到参与的扩大与进入门槛，而不能将立宪语言直接等同于民主实践。

\subsection{全球交换与财政社会}

甲午战后赔款、借款、铁路、租借地、日俄战争和国际贸易，将地方财政同全球资本、商品和军事竞争接合。条约和借款合同能确证法律承诺，不能替代铁路沿线、矿区、港口和农村的实际经验。外资、技术与教育项目往往需要地方土地、劳力、税收和翻译中介，既不能被归结为单向侵略，也不应写成无冲突的发展。

上海、天津、汉口、广州、福州等口岸的银行、轮船、工厂、报馆和学校与省内腹地存在联系，却不构成一条均质的现代城市走廊。价格、工资、租赁和就业的材料多保存于企业、租界或政府文书，带有机构偏向；不能由某一公司的盈利或一座城市的电灯推断全省生活水准。

\subsection{新军、铁路与革命的不同时间}

一九〇八年至一九一一年，新军、铁路国有化争议、保路运动、地方议会与革命组织在不同省份相互牵动。武昌起义及各省响应的日期可由电报、报刊和档案追踪，消息如何抵达县镇、谁控制武器和税粮、家庭怎样面对治安与债务，具有不同时间表。革命不是一夜间抹去旧制度的事件，旧有官吏、商会、宗族、学校和宗教网络仍参与秩序重组。

一九一二年帝制结束为政治制度提供明确终点，却不应成为社会史的句号。改革、教育、立宪、性别关系和全球交换的遗产继续通过省份财政、地方学校、家庭策略和市场道路发挥作用。本节以来源范围限制断言：国家档案最清楚地显示其自身行动，地方材料最能照亮局部实践；两者之间的沉默不能由线性“近代化”叙事替代。
